\section{Primitive events and a valuation obstruction}\label{sec:events}

Let $\Omega=\{0,1\}^{\cS}$ be the cube of bit assignments for
Theorem~\ref{thm:bit-test-system}. A predicate is a Boolean function
$R:\Omega\to\{0,1\}$. We identify predicates that are equal as
functions, regardless of their sampled descriptions. We write
$R\subseteq R'$ when $R(y)\leq R'(y)$ for every $y\in\Omega$.

A \emph{primitive type} is a probability law $P_s$ on predicates,
together with a specified number $p_s\in[0,1]$. An \emph{ordered
comparison} is a coupling of two specified primitive laws supported
on pairs $R\subseteq R'$. The type index is part of the sampling data,
but is not part of the predicate to which a valuation is applied.

\begin{definition}
A Boolean valuation $H$ on the actual predicates is
\emph{$\zeta$-calibrated and ordered} if, for every primitive type
and every specified comparison,
\begin{equation}\label{eq:primitive-calibration-order}
 \abs{\E_{R\sim P_s}H(R)-p_s}\leq\zeta,
 \qquad
 \Prob[H(R)>H(R')]\leq\zeta.
\end{equation}
\end{definition}

\begin{lemma}\label{lem:event-obstruction}
There is a fixed finite family of primitive types $(P_s,p_s)$ and
ordered comparisons, and a fixed $\zeta>0$, constructible in
polynomial time from the input formula, with the following
properties.
\begin{enumerate}[label=\textup{(\roman*)}]
\item Each predicate depends on boundedly many bit positions. All
samplers are polynomially enumerable finite rational laws with
inverse-polynomial positive seed probabilities.
\item Constant-one and constant-zero types are included, with
calibrations one and zero. Every comparison marginal is a primitive
type, and every type has a comparison to the constant-one type.
\item In a satisfiable instance, evaluation at the perfect assignment
$y^*$ satisfies $\E_{P_s}R(y^*)=p_s$ for every type.
\item In an unsatisfiable instance, no Boolean valuation is
$\zeta$-calibrated and ordered.
\end{enumerate}
The numbers $p_s$, the number of types and comparisons, and $\zeta$
are independent of the input size.
\end{lemma}

\subsection{The primitive types}

Use Theorem~\ref{thm:bit-test-system} with pair bound
$\lambda\leq10^{-6}$. There are two basic literal types:
\[
 L_i^a(y)=\ind_{\{y_i=a\}},
 \qquad
 L_{i,j}^a(y)=\ind_{\{y_i+y_j=a\}},
\]
where $i$ is uniform in $\cS$, $(i,j)$ is uniform in $\cS^2$,
and the target $a\in\F_2$ is an independent uniform bit.
Both calibrations are $1/2$.

We use three even-parity tuple experiments:
\begin{enumerate}
\item Four base expressions $y_{i_1},\ldots,y_{i_4}$, with
$(i_1,\ldots,i_4)$ from the auxiliary law of
Theorem~\ref{thm:bit-test-system}.
\item Four sum expressions
$y_{i_1}+y_{j_1},\ldots,y_{i_4}+y_{j_4}$, using two independent
copies of that auxiliary law.
\item The three expressions $y_i,y_j,y_i+y_j$, for independent
uniform $i,j$.
\end{enumerate}
For each experiment of length $k\in\{3,4\}$, independently sample a
uniform even-parity target vector $a\in\F_2^k$. Add the type of the
match event
\[
 M(y)=\ind_{\{(\text{the }k\text{ expressions at }y)=a\}}
\]
and the type of its complement $\overline M=1-M$.
Their calibrations are $2^{-(k-1)}$ and $1-2^{-(k-1)}$.
For each occurrence, add the comparisons
\begin{equation}\label{eq:match-complement-comparisons}
 M\subseteq L^{a_\ell},
 \qquad L^{1-a_\ell}\subseteq\overline M
\end{equation}
with the questions and targets sampled in that experiment.
Their literal marginals are the basic base or sum type, since each
target coordinate is uniform and independent of the questions.
In the perfect assignment all three response tuples have even
parity, so exactly one of the $2^{k-1}$ targets matches.

For each ordinary law in Theorem~\ref{thm:bit-test-system}, let its
arity be $L$ and let $T_e$ be its acceptance predicate for public
seed $e$, with questions $i_1,\ldots,i_L$. Independently choose
$a\in\F_2^L$ uniformly and add the match-and-pass event
\begin{equation}\label{eq:match-and-pass}
 R_{e,a}(y)=T_e(a)\prod_{\ell=1}^L L_{i_\ell}^{a_\ell}(y),
 \qquad p_s=2^{-L}.
\end{equation}
It has the stated calibration in the perfect assignment, which
passes every seeded test. Compare it to each containing literal
$L_{i_\ell}^{a_\ell}$, and include the induced literal marginal as
a type of calibration $1/2$. These additional question marginals
are at most the fixed factor $B$ times uniform.
If the seeded test rejects $a$, Equation~\ref{eq:match-and-pass}
is the constant-zero predicate.

Finally include the two constant types and, for every type $P_s$,
the comparison $(R,1)$ with $R\sim P_s$. This specifies all data in
Lemma~\ref{lem:event-obstruction}. The finite sampling and arity
properties follow from Theorem~\ref{thm:bit-test-system}, since all
products and target spaces have fixed size. We prove its
unsatisfiable-instance assertion next.

\subsection{A quantitative tag--dictator dichotomy}

Suppose $H$ satisfies Equation~\ref{eq:primitive-calibration-order}.
For either the base or sum question space, write
\[
 B(q,a)=2H(L_q^a)-1\in\{-1,1\},\qquad
 b(q)=\frac{B(q,0)+B(q,1)}2.
\]
A question is a \emph{tag} when its two signs agree, in which case
$b(q)=\pm1$; otherwise it is a \emph{dictator}, with $b(q)=0$.
Let $\delta$ be the tag fraction. Literal calibration gives
$|\E b|\leq2\zeta$, and $\E b^2=\delta$.

Consider its four-occurrence experiment and a uniform even target
vector, and put $B_\ell=B(q_\ell,a_\ell)$. Distinct target bits are
independent. The question pair operator has norm at most $\lambda$,
also for sums by the tensor assertion of
Theorem~\ref{thm:bit-test-system}. Hence, for $\ell\ne r$,
\begin{equation}\label{eq:literal-sign-correlation}
 \E B_\ell B_r
 =\E b(q_\ell)b(q_r)
 \leq4\zeta^2+\lambda\delta.
\end{equation}
The first comparisons in
Equation~\ref{eq:match-complement-comparisons} give
\[
 \Prob[B_1=\cdots=B_4=1]\geq\frac18-5\zeta.
\]
For the second comparisons, the event $H(\overline M)=0$ has
probability at least $1/8-\zeta$, and outside the four comparison
failures all opposite literals have value zero. Complementing all
four target bits preserves the uniform even-target law. Therefore
\[
 \Prob[B_1=\cdots=B_4=-1]\geq\frac18-5\zeta.
\]
The squared sum of four signs is $16$ on these two constant-sign
vectors, $4$ when their product is negative, and zero otherwise.
Combining this identity with
Equation~\ref{eq:literal-sign-correlation} gives, for $\zeta\leq1$,
\begin{equation}\label{eq:negative-product-bound}
 \Prob[B_1B_2B_3B_4=-1]\leq
 3\lambda\delta+40\zeta+12\zeta^2
 \leq3\lambda\delta+52\zeta.
\end{equation}

When some but not all questions are tags, the product of the four
signs contains the product of a nonempty proper subset of the target
signs. Under uniform even targets that product is an unbiased sign.
Thus negative product has conditional probability one half.
On the other hand, the probability that the first question is a tag
and the second is not is at least
$\delta(1-\delta)-\lambda\delta$, by the pair-operator bound applied
to the tag indicator. Consequently
\[
 \delta(1-\delta)\leq7\lambda\delta+104\zeta.
\]
For $\lambda\leq1/28$, if $\delta\leq1/2$ this implies
$\delta\leq416\zeta$, and if $\delta>1/2$ it implies
$1-\delta\leq7\lambda+208\zeta$. We have proved
\begin{equation}\label{eq:tag-dichotomy}
 \delta\leq416\zeta
 \quad\text{or}\quad
 1-\delta\leq7\lambda+208\zeta
\end{equation}
for each of the two basic literal types.

\subsection{The tag regime is impossible}

We give explicit error bounds for this step. Assume
$\lambda\leq10^{-6}$ and $\zeta\leq10^{-8}$, and suppose the base
literal type is in the second regime of
Equation~\ref{eq:tag-dichotomy}. Put
\[
 e_0=7\lambda+208\zeta,\qquad
 h_0=e_0+2\zeta,\qquad D_0=3e_0+6\zeta.
\]
Extend the base tag sign arbitrarily to a sign $b_i$ on all of
$\cS$. The non-tag fraction is at most $e_0$, so
$|\E b_i|\leq h_0$. If $r_s=\Prob[b_i=s]$ for $s=\pm1$, then
$|r_s-1/2|\leq h_0/2$ and $r_s^2\leq1/4+h_0$.

In the triple experiment, match calibration and its three
containing-literal comparisons give an all-positive literal-sign
vector with probability at least $1/4-4\zeta$ under even targets.
Complement calibration and the opposite-literal comparisons give
an all-negative vector with the same lower bound under the
complemented, now odd, target law. Outside the event that one of
the two base questions is not a tag, of probability at most $2e_0$,
these occurrences require $b_i=b_j=1$ and $b_i=b_j=-1$,
respectively.
The sum target is a uniform bit under both target laws, independent
of $i,j$. Subtracting these lower bounds from the respective masses
$r_s^2$ yields
\begin{equation}\label{eq:same-class-sum-error}
 \Prob_{i,j,a}\bigl[b_i=b_j=s,\ B((i,j),a)\ne s\bigr]
 \leq h_0+4\zeta+2e_0=D_0,
 \qquad s=\pm1.
\end{equation}
If a sum question is not a tag of sign $s$, at least one of its two
targets has sign different from $s$. Thus
\[
 \Prob_{i,j}[b_i=b_j=s,\ (i,j)\text{ is not a tag of sign }s]
 \leq2D_0.
\]
Since $r_+^2+r_-^2\geq1/2$, the sum tag fraction is at least
$1/2-4D_0>416\zeta$. Its dichotomy in
Equation~\ref{eq:tag-dichotomy} therefore also gives non-tag
fraction at most $e_0$.

Let $v(i,j)$ be the sum tag sign, extended to the remaining pairs
by a fixed sign. We take a symmetric extension: the literal
predicates for $(i,j)$ and $(j,i)$ are equal, so the actual tag
statuses and tag signs are already symmetric. In particular
\begin{equation}\label{eq:sum-sign-on-classes}
 \Prob[b_i=b_j=s,\ v(i,j)\ne s]\leq2D_0.
\end{equation}
For two independent auxiliary four-tuples $\boldsymbol i$ and
$\boldsymbol j$, Equation~\ref{eq:negative-product-bound} and the
at-most $4e_0$ chance of a non-tag sum question imply
\begin{equation}\label{eq:sum-four-product}
 \Prob\left[\prod_{\ell=1}^4v(i_\ell,j_\ell)=-1\right]
 \leq E_4:=3\lambda+52\zeta+4e_0.
\end{equation}
Exchangeability permits swapping $j_1$ and $j_2$ without changing
this bound. Where both fourfold products are positive, cancellation
of their last two factors gives
\[
 v(i_1,j_1)v(i_2,j_2)
 =v(i_1,j_2)v(i_2,j_1).
\]
Thus this rectangle identity fails with probability at most $2E_4$
when $(i_1,i_2)$ and $(j_1,j_2)$ are independent copies of the
auxiliary pair marginal.

We next replace both pair marginals by independent uniform draws.
For fixed $i,i'$, set $f(j)=v(i,j)v(i',j)$. The pair-operator bound
gives
\[
 \abs{\E_{\rm pair}f(j)f(j')-(\E f)^2}
 \leq\lambda\Var(f)\leq\lambda.
\]
This changes the expectation of the rectangle sign by at most
$\lambda$. After replacing the $j$ pair, apply the same calculation
to the $i$ pair with $j,j'$ fixed. The rectangle sign has therefore
changed in expectation by at most $2\lambda$. Its failure
probability for four independent uniform indices is at most
\begin{equation}\label{eq:independent-rectangle-error}
 R_0:=2E_4+\lambda.
\end{equation}

Average over an independent reference pair $(i_0,j_0)$ and fix one
for which the rectangle failure probability over $i,j$ is at most
$R_0$. Define
\[
 u(i)=v(i,j_0),\qquad
 w(j)=v(i_0,j)v(i_0,j_0).
\]
Then
$\Prob[v(i,j)\ne u(i)w(j)]\leq R_0$.
Symmetry of $v$ gives
$\Prob[u(i)w(j)\ne u(j)w(i)]\leq2R_0$.
Equivalently, the signs $a(i)=u(i)w(i)$ at two independent indices
disagree with probability at most $2R_0$.
Choose the majority sign $c$ of $a$. If its minority mass is
$m\leq1/2$, the disagreement probability is $2m(1-m)\geq m$,
so $m\leq2R_0$. Therefore $w(i)=cu(i)$ outside a set of mass at
most $2R_0$, and
\begin{equation}\label{eq:symmetric-sign-factorization}
 \Prob[v(i,j)\ne c\,u(i)u(j)]\leq3R_0.
\end{equation}

Take the base-sign class $s=-c$. Since $h_0<1/2$, it has mass
$r_s\geq1/4$. On two independent draws from this class, the
probability that $u(i)u(j)=1$ is at least one half: if the conditional
positive mass of $u$ is $a$, this probability is
$a^2+(1-a)^2\geq1/2$.
Consequently
\[
 \Prob[b_i=b_j=s,\ c\,u(i)u(j)\ne s]\geq r_s^2/2\geq1/32.
\]
Equations~\ref{eq:sum-sign-on-classes} and
\ref{eq:symmetric-sign-factorization} upper-bound the same
probability by $2D_0+3R_0$. But direct substitution gives
\[
 2D_0+3R_0=231\lambda+6564\zeta<1/32
\]
under the stipulated bounds. The other inequalities used above,
$h_0<1/2$ and $1/2-4D_0>416\zeta$, follow from the same bounds.
This contradiction excludes the base-tag regime. We conclude
\begin{equation}\label{eq:base-dictators}
 \Prob_i[i\text{ is not a dictator}]\leq416\zeta.
\end{equation}

\subsection{Decoding the ordinary tests}

At each dictator question $i$, assign the unique bit $y_i$ whose
literal has $H(L_i^{y_i})=1$; assign arbitrary bits to the other
questions. This is one assignment on $\cS$, because $H$ is a
function of the actual predicates.
For an ordinary law of arity $L$, its probability of querying any
non-dictator is at most $416LB\zeta$ by
Equation~\ref{eq:base-dictators} and the marginal bound.
Its $L$ containing-literal comparisons fail with total probability
at most $L\zeta$ in the seed-and-target experiment.
The constant-zero type has $H(0)=0$, since $\zeta<1$.

Outside these two bad events, if $H(R_{e,a})=1$ in
Equation~\ref{eq:match-and-pass}, all targets agree with the assigned
bits and $T_e(a)=1$. Conversely, an independent uniform target
equals the assigned response vector with probability exactly
$2^{-L}$, even with repeated questions. It follows that
\[
 2^{-L}-\zeta
 \leq \E H(R_{e,a})
 \leq 2^{-L}\Prob_e[T_e(y_{i_1},\ldots,y_{i_L})=1]
       +(416LB+L)\zeta.
\]
Thus the failure probability of this law is at most
\[
 2^L(1+L+416LB)\zeta.
\]
There are only finitely many ordinary laws. Choose the fixed
$\zeta\leq10^{-8}$ small enough that this expression is at most
$\eta/2$ for each of them, where $\eta$ is the soundness constant
of Theorem~\ref{thm:bit-test-system}. In an unsatisfiable instance
this contradicts that theorem. This completes the proof of
Lemma~\ref{lem:event-obstruction}.

For the next section, fix any positive mixture weights
$(\omega_s)_s$ with $\sum_s\omega_s=1$, for example equal weights,
and let $P_*=\sum_s\omega_sP_s$, retaining the type as sampling
data. Put
\begin{equation}\label{eq:primitive-mixture}
 p_*=\sum_s\omega_sp_s.
\end{equation}
The two constant types imply $0<p_*<1$. All these quantities are
fixed independently of the instance size. Equal predicates in
different types remain the same argument of a valuation; only
their sampling probabilities differ.
